\section{Discussion}
\label{sec:discussion}

\para{What do the results show?} The data support a narrow but useful claim: induction-like heads in \gptsmall do affect repeated-token bias on prompts that ask for a random answer. The data do \emph{not} support the stronger claim that induction heads are a large singular source of leakage. The best evidence against that stronger claim is the sensitivity analysis. Depending on which high-scoring heads we ablate, repeated-token bias can either fall or rise. That is hard to reconcile with a monolithic circuit explanation.

\para{Why does the final run still matter?} The final run shows a statistically strong reduction in repeated-token bias, and that effect survives comparison to the random-head control. If we only looked at that run, we might conclude that induction heads are a promising mitigation target. The problem is that the same intervention reduces entropy and harms WikiText utility, while nearby head subsets reverse the direction of the bias effect. The result is therefore real but not robust enough to justify a clean mechanistic story.

\para{How should we interpret the entropy drop?} Lower repeated-token bias is not the same as better randomness. In our setup, ablating \headpair reduces the probability assigned to the repeated distractor, but it also makes the candidate distribution less diffuse overall. That pattern echoes the broader randomness-control literature \cite{zhang2024diffuse}: moving one undesired mode down does not guarantee a more uniform distribution across the valid options. Mechanistic mitigation needs to track distributional quality directly, not only the probability of one repeated token.

\para{Connection to memorization and leakage work.} Our results fit better with the literature that treats memorization and leakage as distributed phenomena \cite{lee2022deduplicating, ippolito2022falseprivacy, huang2024memorization}. Hard ablation can change one observable behavior, but the effect is entangled with general model quality and with the interaction among heads. That does not make mechanistic analysis unhelpful. It means the analysis should aim for richer behavior-specific interventions, such as softer descaling \cite{wang2025toxicity}, rather than assuming one clean deletion will isolate leakage.

\para{Limitations.} This study tests only one model family, \gptsmall. The random-choice task is synthetic and uses a base model rather than an instruction-tuned model, so the prompt instruction to choose randomly is weaker than it would be in a chat model. CPU-only execution limited the final sample sizes and the breadth of sensitivity sweeps. The utility check uses only 12 WikiText samples, which is enough to show a large degradation but not enough to characterize subtle language modeling effects. Finally, we do not run a true training-data extraction attack, so our ``leakage'' metric should be read as prompt-induced copying pressure rather than direct memorization extraction.

\para{Broader implications.} The main practical lesson is methodological. Circuit-level interventions for leakage should not be evaluated on a single synthetic metric in isolation. They need at least three checks: whether the target behavior changes, whether the output distribution remains well-formed, and whether general utility survives. Our results show that a promising-looking mechanistic effect can fail that broader evaluation.
